\section{Discussion and limitations}\label{sec:discussion}
In the simulation family~\eqref{eq:family}, the error of a regression GATE estimator is a variance that follows the corner leverage and a bias that equals the exact projection bias of the response. The variance needs only the graph and $p$; the bias needs the shape of the response, which an experiment cannot supply.
Three hypotheses we held did not survive: there is no crossover with exact Horvitz--Thompson at a moderate degree; an error in the observed graph is not graceful (10\% wrong edges raise the error by half); and three of the four data-only scores did not beat the corner distance as detectors of bias.
For use: regression is the better estimator over most of the map (Table~\ref{tab:map}); for a suspected threshold at $p=\tfrac12$, compare with a relaxed exposure estimator; a design that puts the corner in the data removes much of the bias at some cost.

\textbf{Limitations.} All results are for one parametric family with synthetic outcomes; the real graphs supply topology only, so they test degree heterogeneity, not whether real interference follows this family. Sizes are small ($n\approx2000$, up to $16{,}000$), most grids use Bernoulli assignment, and the diagnostics are evaluated at one $n$.
